\section{Results}
\label{sec:results}
The following results distinguish published WORKS endpoints, retained
component measurements, and the new live transfer executions.

\subsection{Memory-mapped file accesses and feature contributions}
Table~\ref{tab:mmap} compares explicit reads alone with reads plus
memory-mapped file accesses, as briefly reported in
WORKS~\cite{camp2026works}. Adding the mapped-access signal reduces
underallocations from 32 to 30 and total requests by 18.11\%.
Median prediction error increases from 22.04\% to 26.51\%. In this
evaluation, mapped-access measurements therefore yield fewer
underallocations and lower total requests alongside higher median
prediction error.

\begin{table}[!t]
\caption{Memory-mapped access ablation on 5,026 tasks.}
\label{tab:mmap}
\centering\small
\input{\campReadyPath result_assets/mmap_table.tex}
\end{table}

Table~\ref{tab:features} shows that CAMP has the fewest underallocations:
30 versus 85 for Access, 46 with consumption history, and 52 with
predicted consumption. Relative to Access, CAMP reduces unused memory-time
by 37.58\% while requesting 31.43\% more memory; MdAPE increases from
21.75\% to 26.51\%.

Among consumption-enhanced variants, full CAMP requests the least memory.
Adding predicted consumption to consumption history reduces underallocations
from 46 to 30 and requests by 57.61\%, as summarized in
WORKS~\cite{camp2026works}. Consumption history alone has the lowest MdAPE
(21.32\%) and unused memory-time (37.19~GiB-hours versus CAMP's 39.75).
The fewest underallocations therefore coincide with different settings
from the lowest prediction error and unused memory-time.

\begin{table}[!t]
\caption{Feature ablation on 5,026 tasks.}
\label{tab:features}
\centering\footnotesize
\setlength{\tabcolsep}{3pt}
\input{\campReadyPath result_assets/feature_table.tex}
\end{table}

\subsection{Downstream allocation and auditing cost}
At the primary 20\% budget, selective auditing yields 265.8 mean
underallocations versus 305.8 for random auditing (13.1\% fewer), while
mean requests increase from 53,123.35 to 54,400.10~GiB (2.4\%).
WORKS also summarized this result~\cite{camp2026works}.
Figure~\ref{fig:audit}(a) shows lower mean underallocation rates for
selective auditing at 5\%, 20\%, and 30\%, and higher rates at 10\% and
50\%.
At 20\%, the paired 95\% confidence interval for the underallocation-rate
difference is $[-0.416,0.007]$ percentage points and includes zero.
The mean favors selective auditing at this budget, while the interval
includes differences in either direction.

Figure~\ref{fig:audit}(b) shows estimated auditing cost $C_B$ at budgets from
5\% to 50\%. Mean cost increases from 5.999~GiB-hours at 5\% to
24.025~GiB-hours at 20\% and 66.534~GiB-hours at 50\%. At 20\%, the
95\% confidence interval is $[23.250,24.673]$~GiB-hours. Full auditing
costs 139.680~GiB-hours. These estimates quantify additional reserved
memory-time from auditing under fixed CAMP requests, separately from the
downstream allocation comparison.

\subsection{Repeated inputs and transfer to new inputs}
Figure~\ref{fig:deployment} summarizes the four execution stages.
Cold-to-warm requests fall from 977.79 to 241.12~GiB (75.34\%),
while aggregate peaks remain near 229.3~GiB. Unused memory-time falls
from 32.86089 to 0.01847~GiB-hours. Warm coverage is 99.58\% and
MdAPE is 1.13\%, showing tighter reservations with repeated-input feedback.

\begin{figure*}[!t]
\centering
\includegraphics[width=\textwidth]{\campReadyPath result_assets/online_deployment_compact.pdf}
\caption{Live Bowtie2 deployment, 6,000 tasks per stage. Cold and Warm use original inputs; WT and DT use the same new inputs after warm and cold checkpoints. Panels show (a) peaks and requests, (b) unused memory-time, and (c) request margins. Requests exclude the 128~MiB execution allowance; below-peak requests are not recorded OOM failures.}
\label{fig:deployment}
\end{figure*}

On new inputs, Warm-Transfer peaks rise to 273.44~GiB and requests to
319.22~GiB. Coverage falls to 94.15\%, with 351 requests below peak
versus 25 on repeated inputs (Fig.~\ref{fig:deployment}(c)). Overall
MdAPE rises from 1.13\% to 4.73\%; alignment MdAPE rises from 0.91\%
to 15.60\%. Alignment accounts for 1,974 of 6,000 tasks and exhibits
a larger error increase than the overall population.

\subsection{Effect of the extra warm pass}
On the same new inputs, Direct-Transfer requests 343.56~GiB, 7.62\%
more than Warm-Transfer, with unused memory-time of 0.22109 versus
0.12805~GiB-hours (72.65\% higher). Coverage improves from 94.15\%
to 95.78\%, with 253 rather than 351 requests below peak.
Warm-Transfer therefore reduces reservations and unused memory-time,
while Direct-Transfer achieves higher coverage.
Overall MdAPE favors Direct-Transfer (2.69\% versus 4.73\%);
alignment MdAPE favors Warm-Transfer (15.60\% versus 18.96\%).
Peak totals differ by 1.65\% between branches.

Each stage completes 6,000 tasks with zero recorded OOMs under the
execution allowance in Fig.~\ref{fig:deployment}. Cold, warm, WT, and
DT record 9, 15, 5, and 10 retries. Both new-input branches produce
identical flagstat summaries for all 1,974 chunks. The reported $W$
uses successful-attempt runtimes, excluding failed attempts; auditing
and learning costs are not separately accounted for.
